For the finite quadrupole condition, consider the general map
\[
 V=\Sym_0(n,\mathbb R),\qquad
 \mu:V\times V\longrightarrow\so(n),\qquad \mu(A,B)=[A,B],
 \qquad Z=\mu^{-1}(0),\qquad n\geq2.
\]
The metrics are the Frobenius metric on matrices and its product metric on
$V\times V$.  This is standard commuting symmetric-matrix geometry.  Under
the canonical identification $T^*V\simeq V\times V$, the commutator represents
the moment map for simultaneous orthogonal conjugation, up to the sign fixed
by the dual-pairing convention.  The claims below do not concern the range
of a gravitational constraint operator, interface conditions, or nonlinear
spacetime continuation.  The momentum-map context is classical \cite{ArmsMarsdenMoncrief1981}.

Fix $(A,B)\in Z$.  Simultaneous orthogonal diagonalization gives the
\emph{maximal} joint eigenspace decomposition
\[
 \mathbb R^n=\bigoplus_{s=1}^r E_s,\qquad
 A|_{E_s}=\alpha_s I,\qquad B|_{E_s}=\beta_s I,
 \qquad (\alpha_s,\beta_s)\ne(\alpha_t,\beta_t)\quad(s\ne t).
\]
Write $m_s=\dim E_s$, let $P_s$ be its orthogonal projector, and put
$a_{st}=P_s a|_{E_t}$ and $a_{\mathrm d}=\sum_sP_saP_s$, with analogous
notation for $b$.  In a simultaneous orthonormal eigenbasis, use repeated
lists $\alpha_i,\beta_i$ and define
\[
 c_{ij}=\beta_j-\beta_i,\qquad d_{ij}=\alpha_i-\alpha_j,
 \qquad g_{ij}^2=c_{ij}^2+d_{ij}^2\quad(i<j).
\]

\begin{proposition}[Rank, optimal inverse, and regularity]\label{prop:inverse}
The derivative $L=D\mu_{(A,B)}$ is
\[
 L(a,b)=[a,B]+[A,b],\qquad L(a,b)_{ij}=c_{ij}a_{ij}+d_{ij}b_{ij}.
\]
Its nonzero singular values are exactly the positive $g_{ij}$.  Equivalently,
for each $s<t$ the singular value
\[
 \sqrt{(\alpha_s-\alpha_t)^2+(\beta_s-\beta_t)^2}
\]
occurs with multiplicity $m_s m_t$.  Hence, with
$d=\dim V=n(n+1)/2-1$ and $N=\dim\so(n)=n(n-1)/2$,
\[
 \operatorname{rank} L=N-\sum_s\binom{m_s}{2},\qquad
 \dim\ker L=2d-N+\sum_s\binom{m_s}{2}.
\]
The range comprises precisely the skew matrices $S$ with $P_sSP_s=0$ for
all $s$.  The Moore--Penrose inverse $R=L^\dagger$ is the pair of symmetric
zero-diagonal matrices with
\[
 (RS)^{(1)}_{ij}=\frac{c_{ij}S_{ij}}{g_{ij}^2},\qquad
 (RS)^{(2)}_{ij}=\frac{d_{ij}S_{ij}}{g_{ij}^2}\qquad(g_{ij}>0),
\]
and both entries zero when $g_{ij}=0$.  Here the two displayed components
are extended symmetrically.  Then $LR$ is the orthogonal range projection.
If at least one $g_{ij}>0$,
\[
 \|R\|_{\so(n)\to V\times V}
   =\frac{1}{\min_{g_{ij}>0}g_{ij}}.
\]
For $L=0$, $R=0$ and $\|R\|=0$.  In particular, a right inverse on all of
$\so(n)$ exists exactly when every joint eigenvalue is simple.  In that
case $R$ is the unique solution of minimum product norm for each target,
and its operator norm is the smallest possible among all linear right
inverses.
\end{proposition}

\begin{proof}
For $i<j$, the normalized symmetric and skew matrix units are respectively
$(E_{ij}+E_{ji})/\sqrt2$ and $(E_{ij}-E_{ji})/\sqrt2$.  On the two normalized
input directions belonging to this edge, $L$ is the single row
$(c_{ij},d_{ij})$.  Distinct edges are orthogonal; all diagonal input
directions are in the kernel.  The trace restrictions affect only those
diagonal inputs.  Thus each edge contributes one singular value $g_{ij}$
unless it lies within a joint eigenspace.  This proves the singular-value,
rank, and range statements.  The inverse formula is the minimum-norm
solution of the scalar equation $c_{ij}a_{ij}+d_{ij}b_{ij}=S_{ij}$.
Consequently,
\[
 \|RS\|_{V\times V}^2
   =2\sum_{i<j:\,g_{ij}>0}\frac{S_{ij}^2}{g_{ij}^2}.
\]
The norm formula follows by concentrating $S$ on an edge of minimum
positive gap.  Every solution for that target has norm at least the
minimum-norm solution, which also proves optimality among right inverses.
There is no extra factor $\sqrt2$ in the singular values: the input and
output Frobenius normalizations agree.
\end{proof}

\begin{theorem}[All velocities of commuting curves]\label{thm:velocities}
For $a,b\in V$, the following conditions are equivalent.
\begin{enumerate}
\item There is a $C^1$ curve $t\mapsto(A(t),B(t))\in Z$ through $(A,B)$
with derivative $(a,b)$ at $0$.
\item $L(a,b)=0$ and $[a_{ss},b_{ss}]=0$ on every maximal joint eigenspace
$E_s$.
\item $L(a,b)=0$ and there are unrestricted $c,d\in V$ satisfying the
second-order matrix equation
\[
 L(c,d)=-2[a,b].
\]
\item There is a real-analytic such curve with derivative $(a,b)$ at $0$.
\end{enumerate}
For a velocity satisfying these conditions, define a skew matrix $K$ by
$K_{ss}=0$ and, for $s\ne t$,
\[
 K_{st}=
 \frac{(\alpha_t-\alpha_s)a_{st}+(\beta_t-\beta_s)b_{st}}
 { (\alpha_t-\alpha_s)^2+(\beta_t-\beta_s)^2 }.
\]
One explicit realization is
\begin{align*}
 Q(t)&=(I+tK/2)(I-tK/2)^{-1},\\
 A(t)&=Q(t)(A+t a_{\mathrm d})Q(t)^{\mathsf T},&
 B(t)&=Q(t)(B+t b_{\mathrm d})Q(t)^{\mathsf T}.
\end{align*}
This rational curve is defined and real analytic for every real $t$.
\end{theorem}

\begin{proof}
First, $L(a,b)=0$ is necessary by differentiation.  Between distinct joint
blocks it reads
\[
 (\beta_t-\beta_s)a_{st}-(\alpha_t-\alpha_s)b_{st}=0.
\]
The coefficient pair is nonzero.  Therefore the displayed definition of
$K$ gives
\[
 a_{st}=(\alpha_t-\alpha_s)K_{st},\qquad
 b_{st}=(\beta_t-\beta_s)K_{st}\quad(s\ne t).
\]
Symmetry of $a,b$ and reversal of the gaps imply $K_{ts}=-K_{st}^{\mathsf T}$.
In particular,
\begin{equation}\label{eq:offblock-cancellation}
 a_{st}b_{ts}-b_{st}a_{ts}=0\quad(s\ne t),\qquad
 P_s[a,b]P_s=[a_{ss},b_{ss}].
\end{equation}

For necessity of the block condition, differentiability at $0$ alone is
enough.  Put $X(t)=A(t)-A=ta+o(t)$ and $Y(t)=B(t)-B=tb+o(t)$.  Every
commutator with $A$ or $B$ has zero diagonal joint block, so the exact
equation $[A+X(t),B+Y(t)]=0$ gives
\[
 0=P_s[X(t),Y(t)]P_s=t^2P_s[a,b]P_s+o(t^2).
\]
Divide by $t^2$ and use \eqref{eq:offblock-cancellation}.  This proves
$(1)\Rightarrow(2)$ without assuming a second derivative of the curve.

For sufficiency, the block condition says $[a_{\mathrm d},b_{\mathrm d}]=0$.
Both block-diagonal matrices commute with $A$ and $B$, since the latter
are scalar on each joint block.  Thus
$[A+t a_{\mathrm d},B+t b_{\mathrm d}]=0$ identically.  For real $t$,
$I-tK/2$ is invertible: if $(I-tK/2)x=0$, taking the real inner product with
$x$ gives $\|x\|^2=0$.  The Cayley matrix $Q(t)$ is orthogonal, has
$Q(0)=I$, and has $Q'(0)=K$.  Hence the proposed curve is symmetric,
tracefree, real analytic, and commuting, and
\[
 A'(0)=a_{\mathrm d}+[K,A]=a,\qquad
 B'(0)=b_{\mathrm d}+[K,B]=b.
\]
This proves $(2)\Rightarrow(4)\Rightarrow(1)$.

Finally, Proposition~\ref{prop:inverse} says the second-order equation is
solvable exactly when $P_s[a,b]P_s=0$ for all $s$.  Under $L(a,b)=0$,
\eqref{eq:offblock-cancellation} identifies this with the block condition.
This proves $(2)\Leftrightarrow(3)$.  Explicit accelerations of the curve are
\[
 c=2[K,a_{\mathrm d}]+[K,[K,A]],\qquad
 d=2[K,b_{\mathrm d}]+[K,[K,B]],
\]
and $R(-2[a,b])$ is the correction of minimum product norm.
\end{proof}

\begin{remark}
The condition is \emph{not} $[a,b]=0$ on all of $\mathbb R^n$.  Off-block
components of this commutator can be canceled by second-order corrections.
For example, let $A=\diag(1,-1,0,0)$, $B=0$,
$a=E_{12}+E_{21}$, and $b=A$.  Then $L(a,b)=0$ and every diagonal joint
block commutator is zero, although $[a,b]\ne0$.  Splitting a repeated
maximal joint block into artificial one-dimensional blocks would instead
discard genuine obstructions and is invalid.
\end{remark}

\begin{corollary}[Smooth points and a simple pencil]\label{cor:regular}
The following are equivalent at a commuting pair: $L$ is surjective; all
maximal joint eigenspaces are one dimensional; some real pencil $A+\tau B$
has simple spectrum; and $Z$ is a $C^1$ embedded submanifold in a
neighborhood of that pair.  Its dimension there is
\[
 2\dim V-\dim\so(n)=\frac{n^2+3n-4}{2}.
\]
Neither $A$ nor $B$ need individually have simple spectrum.
\end{corollary}

\begin{proof}
Surjectivity is Proposition~\ref{prop:inverse}, and the implicit-function
theorem gives a smooth submanifold at every such point.  If $m_s\geq2$,
choose an internal tracefree diagonal matrix $a$ with entries $1,-1$ and
$b=E_{ij}+E_{ji}$ on those two coordinates.  The velocities $(a,0)$ and
$(0,b)$ are realized by straight commuting curves, while their sum fails
Theorem~\ref{thm:velocities}.  The velocity set of a $C^1$ embedded
submanifold would be a vector space.  Thus no repeated joint block is a
smooth point.  For distinct joint pairs, a collision
$\alpha_i+\tau\beta_i=\alpha_j+\tau\beta_j$ excludes at most one value of
$\tau$ per pair; if the two $\beta$ values agree, no value is excluded.
Avoiding this finite set gives a simple pencil.  A repeated joint block
prevents a simple pencil for every $\tau$.
\end{proof}

\noindent
A local polynomial chart at any regular pair can therefore be based on a
simple pencil rather than on $A$ alone.  Fix a suitable $\tau$ and set
$C=A+\tau B$, $D=B$.  Every symmetric tracefree $D$ commuting with a
nearby simple $C$ is uniquely
\[
 D=\sum_{k=1}^{n-1}c_k\left(C^k-\frac{\tr(C^k)}{n}I\right),\qquad
 A=C-\tau D,\quad B=D.
\]
Interpolation on the $n$ distinct eigenvalues proves existence and
uniqueness.  The coefficients depend smoothly on $(C,D)$, for example by
inverting the nonsingular Frobenius Gram matrix of the displayed
polynomial basis.  The $\dim V+n-1$ parameters agree with
Corollary~\ref{cor:regular}.  In dimension four, an example with neither
individual matrix simple is
\[
 A=\diag(-1,-1,1,1),\qquad B=\diag(-1,1,-1,1).
\]
All four joint pairs are distinct, and $A+2B=\diag(-3,1,-1,3)$ is simple.

\paragraph{Dimension-four strata and the audited seeds.}
For $n=4$, the complete multiplicity list gives
\[
\begin{array}{c|r|r}
\text{joint multiplicities}&\operatorname{rank}L&\dim\ker L\\ \hline
1,1,1,1&6&12\\
2,1,1&5&13\\
2,2&4&14\\
3,1&3&15\\
4&0&18
\end{array}
\]
Every row other than the first has a nonlinear set of realizable velocities.
At the norm-two seed $A=B=\diag(1,-1,0,0)$, take
$a=\diag(0,0,1,-1)$ and $b=E_{34}+E_{43}$.  Their sum satisfies the linear
equation but has $[a,b]_{34}=2$ in the repeated block, so no $C^1$ commuting
curve realizes that velocity.  This proves the seed's genuine
singularity, rather than merely a failure of the chosen equations to have
independent derivatives.

For the norm-matched replacement
\[
 A=B=\frac1{\sqrt{10}}\diag(3,1,-1,-3),\qquad
 \tr(A^2)=\tr(B^2)=2,
\]
the minimal single-matrix spectral gap is $\gamma=2/\sqrt{10}$ and the
minimal joint gap is $\sqrt2\gamma=2/\sqrt5$.  Thus the inverse bound is sharp:
\[
 \|R\|=\frac{\sqrt5}{2}.
\]
When $A=B$ has distinct eigenvalues $\lambda_i$, the general inverse
reduces to $a=-C/2$, $b=C/2$, where
$C_{ij}=S_{ij}/(\lambda_i-\lambda_j)$ and $C_{ii}=0$.

\paragraph{Scope for the quadrupole application.}
When the restricted charge equations have independently been established
to reduce to $[A,B]=0$ with a nonzero prefactor, the preceding statements
classify their entire finite curve-velocity obstruction and the optimal
linear inverse.  They make no claim that charge-zero data, or finite
second-order solvability, suffice for the Einstein--scalar bulk and shell
equations.  The linear map $D\mu$ and its range must not be substituted for
the full constraint/interface operator.
